\section{Rendering --- current state}
\label{sec:rendering-current}

\begin{statusbox}
\textbf{Half-migrated.} Interfaces for a command/token backend exist
(\symbolref{IBackend}, \symbolref{RenderCommand}, \symbolref{CommandQueue},
\symbolref{Renderer::enroll*}), but the live frame path still largely uses
\textbf{immediate OpenGL} inside drawables and \symbolref{Scene::Update}.
\end{statusbox}

\subsection{What actually draws frames today}
\begin{enumerate}
  \item Modules call \symbolref{DispatchDrawables} and add \symbolref{DrawableBase*}
        pointers to \symbolref{Scene}.
  \item \symbolref{Scene::Update} clears color/depth with raw \texttt{glClear*},
        then calls \symbolref{DrawableBase::Draw} on each entry.
  \item Concrete drawables (e.g.\ \symbolref{Canvas} via CRTP
        \symbolref{Drawable<Canvas>}) run \symbolref{DrawImpl}, which issues GL
        calls and owns GL object handles (VAO/VBO/textures/programs) on the
        drawable base.
\end{enumerate}

\subsection{Drawable design (observed)}
\begin{itemize}
  \item \symbolref{DrawableBase} --- type-erased base for storage; virtual
        \symbolref{Draw}; destructor deletes GL objects.
  \item \symbolref{Drawable<Derived>} --- CRTP: virtual \symbolref{Draw}
        dispatches to \symbolref{Derived::DrawImpl} for static polymorphism
        after the first virtual hop.
\end{itemize}

\begin{inferencebox}
CRTP here is a micro-optimization / style choice: keep a homogeneous list of
bases while avoiding a second virtual call into each concrete drawer. Resource
ownership on the base tightly couples ``something that can be listed in a scene''
to ``OpenGL object lifetime,'' which fights multi-API abstraction.
\end{inferencebox}

\subsection{Parallel / unfinished stacks still in tree}
\begin{table}[h]
\centering
\begin{tabularx}{\linewidth}{@{}l X@{}}
\toprule
\textbf{Stack} & \textbf{Notes} \\
\midrule
\symbolref{RenderQueue} on window & Older ``queue of drawables'' idea; add paths often commented out. \\
\symbolref{IBackend} + \symbolref{CommandQueue} & Target direction; \symbolref{Renderer::RenderScene} mostly clears the queue and returns. \\
\symbolref{EntityTable} / \symbolref{ModuleObject} & Handle-ish objects with mesh/texture IDs; not the primary cell path. \\
\symbolref{SceneObject} hierarchy & Parent/child transforms; lightly used relative to drawable list. \\
\bottomrule
\end{tabularx}
\end{table}

\subsection{Backend surface (already sketched)}
\symbolref{IBackend} exposes:
\begin{itemize}
  \item lifecycle: \symbolref{Initialize}, \symbolref{Shutdown},
        \symbolref{BeginFrame}, \symbolref{EndFrame}
  \item queue: \symbolref{PushToCommandQueue}, \symbolref{SubmitCommandQueue},
        \symbolref{ClearCommandQueue}
  \item resource creation returning \texttt{unsigned long} handles:
        mesh / shader / texture / descriptor set
\end{itemize}

\symbolref{CommandType} already enumerates a useful token vocabulary:
pipeline state, bindings, uniforms, render-pass ops, draw ops
(\pathref{Rendering/RenderCommand.h}).

\symbolref{Renderer} separates:
\begin{itemize}
  \item \textbf{Enrollment} --- create GPU resources up front (not mixed into
        per-frame command spam).
  \item \textbf{RenderScene} --- intended to translate a scene into tokens and
        submit (stub today).
\end{itemize}

\begin{decisionbox}[title={Intent already visible in code}]
Game/UI code should not call \texttt{gl*} directly long-term. It should either
enroll resources or emit \symbolref{RenderCommand} tokens; only the backend
implementation speaks the API.
\end{decisionbox}

\subsection{Pain points blocking the migration}
\begin{enumerate}
  \item \symbolref{Canvas} is both domain grid and GL texture/drawable.
  \item \symbolref{DrawableBase} owns GL handles.
  \item \symbolref{Scene::Update} is a mini immediate renderer.
  \item \symbolref{RenderCommand} payload is a flat ``kitchen sink'' struct
        (not a tagged union / \texttt{std::variant}); many fields unused per
        command type.
  \item Two drawable submission ideas (\symbolref{Scene} list vs
        \symbolref{RenderQueue}) without a single winner.
\end{enumerate}
